\subsection{Ablation Study}
\label{app:ablation}

To isolate the contribution of each component of Red-LGCP, we compare the full model with six variants, each removing or replacing one component while keeping everything else fixed (Table~\ref{tab:results-ablation}). The first three act on the gate and reuse the trained models of the full model; the other three change the model or its training. Each variant is compared with the full model by a paired two-sided exact Wilcoxon signed-rank test over the eight test weeks, with metrics averaged over seeds as in App.~\ref{app:significance}. These p-values are not corrected for multiple comparisons, but a difference in all eight weeks ($p=0.008$) remains significant after Holm correction across the six variants.

\textbf{w/o uncertainty} uses the fixed threshold $\tau$ instead of $\tau+\sigma_{j,t}$. Requiring a margin of one standard deviation makes the gate more conservative on uncertain cells, which improves accuracy in seven of the eight weeks ($p=0.039$) at the cost of a slightly higher RMSE, lower with the fixed threshold in seven weeks ($p=0.016$); MAE, Wasserstein distance and F1 change little and in no consistent direction.

\textbf{w/o redistribution} zeroes the rejected cells without recovering any of their expected counts. The gate decisions are unchanged, and so are accuracy and F1, but the Wasserstein distance is higher in all eight weeks (0.126 vs.\ 0.101 in May, 0.115 vs.\ 0.085 in November; $p=0.008$), confirming that redistribution restores the overall distribution of predicted counts. MAE is instead slightly lower without it (0.275 in both months; six of eight weeks, $p=0.039$): redistribution trades a small pointwise cost for a better distributional fit.

\textbf{w/o gate} reports the expected counts $\mu_{j,t}$ of the full model without zeroing. These are positive in every cell-day, so accuracy and F1 drop to the level of the covariate-based baselines in Table~\ref{tab:results-main}, and MAE and Wasserstein distance are higher in all eight weeks ($p=0.008$). RMSE is instead slightly lower without the gate (0.903 vs.\ 0.941 in May, 0.849 vs.\ 0.881 in November; six of eight weeks, $p=0.078$). This reflects the two error measures: RMSE is minimized by the mean of the predictive distribution, i.e., the ungated expected count, whereas MAE is minimized by its median, which is zero whenever a cell is more likely inactive than active.

\textbf{two-stage} is the model without joint optimization: the LGCP is fitted alone, with a Poisson likelihood on all counts, and a classifier with the same inputs and hidden layer, but without dropout, is trained afterwards and used as a post-hoc gate with the fixed threshold $\tau$ and the same redistribution. Its log-likelihood and CRPS are those of the Poisson LGCP, since a post-hoc gate only changes point predictions. Training the two components together through the hurdle likelihood relieves the LGCP of explaining the zeros, and improves the log-likelihood, the CRPS and the Wasserstein distance (0.164 vs.\ 0.101 in May, 0.137 vs.\ 0.085 in November) in all eight weeks ($p=0.008$), and accuracy in seven ($p=0.023$). MAE and RMSE are also lower in both months, though not significantly, and vary much less across weeks and seeds (RMSE standard deviation 0.115 vs.\ 0.482 in May), while F1 is similar.

\textbf{w/o classifier} is the LGCP alone, with a Poisson likelihood and no gate. Its rates are positive in every cell-day, so accuracy and F1 drop as for the ungated model, and MAE, Wasserstein distance, log-likelihood and CRPS are all worse than for the full model in every week ($p=0.008$), confirming that the classifier is the main source of the advantage of Red-LGCP over the baselines.

\textbf{w/o kernel} fixes all kernel variances to $10^{-10}$ (non-trainable), which removes the latent field and reduces the intensity to the covariate effect, $\log\lambda_i=\alpha+\mathbf{x}_i^\top\boldsymbol{\beta}$ (\S\ref{app:hyperparameters} details why this reduction is numerically stable). The classifier does not depend on the latent field, so accuracy and F1 are essentially unchanged. The latent field improves all other metrics in May (e.g., RMSE 1.017 vs.\ 0.941, log-likelihood $-0.466$ vs.\ $-0.441$) and leaves November nearly unchanged, so no difference is significant over the eight weeks, suggesting that the spatiotemporal structure adds signal mainly in May, whereas in November the covariates, which include lagged counts, already capture most of it.

\begin{table*}[t]
\centering
\footnotesize
\setlength{\tabcolsep}{2.5pt}
\begin{tabular}{llccccccc}
\toprule
Variant & Month & MAE & RMSE & Wasserstein & Accuracy & F1 & Log-lik. & CRPS \\
\midrule
\multirow{2}{*}{Red-LGCP (full)} & May & 0.286\,{\scriptsize$\pm$0.041} & 0.941\,{\scriptsize$\pm$0.115} & \textbf{0.101}\,{\scriptsize$\pm$0.047} & \textbf{0.912}\,{\scriptsize$\pm$0.009} & \textbf{0.645}\,{\scriptsize$\pm$0.072} & \textbf{$-$0.441}\,{\scriptsize$\pm$0.061} & \textbf{0.208}\,{\scriptsize$\pm$0.042} \\
 & Nov. & 0.283\,{\scriptsize$\pm$0.058} & 0.881\,{\scriptsize$\pm$0.278} & 0.085\,{\scriptsize$\pm$0.015} & \textbf{0.899}\,{\scriptsize$\pm$0.021} & 0.573\,{\scriptsize$\pm$0.096} & \textbf{$-$0.455}\,{\scriptsize$\pm$0.042} & \textbf{0.209}\,{\scriptsize$\pm$0.043} \\
\midrule
\multirow{2}{*}{w/o uncertainty} & May & 0.289\,{\scriptsize$\pm$0.042} & 0.928\,{\scriptsize$\pm$0.126} & 0.104\,{\scriptsize$\pm$0.045} & 0.903\,{\scriptsize$\pm$0.008} & 0.637\,{\scriptsize$\pm$0.065} & \textbf{$-$0.441}\,{\scriptsize$\pm$0.061} & \textbf{0.208}\,{\scriptsize$\pm$0.042} \\
 & Nov. & 0.281\,{\scriptsize$\pm$0.053} & 0.864\,{\scriptsize$\pm$0.263} & 0.090\,{\scriptsize$\pm$0.016} & 0.895\,{\scriptsize$\pm$0.024} & \textbf{0.590}\,{\scriptsize$\pm$0.089} & \textbf{$-$0.455}\,{\scriptsize$\pm$0.042} & \textbf{0.209}\,{\scriptsize$\pm$0.043} \\
\midrule
\multirow{2}{*}{w/o redistribution} & May & \textbf{0.275}\,{\scriptsize$\pm$0.051} & 0.910\,{\scriptsize$\pm$0.141} & 0.126\,{\scriptsize$\pm$0.053} & \textbf{0.912}\,{\scriptsize$\pm$0.009} & \textbf{0.645}\,{\scriptsize$\pm$0.072} & \textbf{$-$0.441}\,{\scriptsize$\pm$0.061} & \textbf{0.208}\,{\scriptsize$\pm$0.042} \\
 & Nov. & \textbf{0.275}\,{\scriptsize$\pm$0.052} & 0.873\,{\scriptsize$\pm$0.257} & 0.115\,{\scriptsize$\pm$0.014} & \textbf{0.899}\,{\scriptsize$\pm$0.021} & 0.573\,{\scriptsize$\pm$0.096} & \textbf{$-$0.455}\,{\scriptsize$\pm$0.042} & \textbf{0.209}\,{\scriptsize$\pm$0.043} \\
\midrule
\multirow{2}{*}{w/o gate} & May & 0.356\,{\scriptsize$\pm$0.047} & \textbf{0.903}\,{\scriptsize$\pm$0.140} & 0.218\,{\scriptsize$\pm$0.044} & 0.131\,{\scriptsize$\pm$0.030} & 0.231\,{\scriptsize$\pm$0.046} & \textbf{$-$0.441}\,{\scriptsize$\pm$0.061} & \textbf{0.208}\,{\scriptsize$\pm$0.042} \\
 & Nov. & 0.344\,{\scriptsize$\pm$0.045} & \textbf{0.849}\,{\scriptsize$\pm$0.241} & 0.200\,{\scriptsize$\pm$0.012} & 0.122\,{\scriptsize$\pm$0.019} & 0.217\,{\scriptsize$\pm$0.030} & \textbf{$-$0.455}\,{\scriptsize$\pm$0.042} & \textbf{0.209}\,{\scriptsize$\pm$0.043} \\
\midrule
\multirow{2}{*}{two-stage} & May & 0.304\,{\scriptsize$\pm$0.079} & 1.244\,{\scriptsize$\pm$0.482} & 0.164\,{\scriptsize$\pm$0.075} & 0.895\,{\scriptsize$\pm$0.013} & 0.632\,{\scriptsize$\pm$0.068} & $-$0.520\,{\scriptsize$\pm$0.082} & 0.247\,{\scriptsize$\pm$0.068} \\
 & Nov. & 0.286\,{\scriptsize$\pm$0.061} & 0.945\,{\scriptsize$\pm$0.374} & 0.137\,{\scriptsize$\pm$0.029} & 0.891\,{\scriptsize$\pm$0.020} & 0.577\,{\scriptsize$\pm$0.084} & $-$0.525\,{\scriptsize$\pm$0.054} & 0.226\,{\scriptsize$\pm$0.044} \\
\midrule
\multirow{2}{*}{w/o classifier} & May & 0.377\,{\scriptsize$\pm$0.078} & 1.206\,{\scriptsize$\pm$0.456} & 0.256\,{\scriptsize$\pm$0.063} & 0.131\,{\scriptsize$\pm$0.030} & 0.231\,{\scriptsize$\pm$0.046} & $-$0.520\,{\scriptsize$\pm$0.082} & 0.247\,{\scriptsize$\pm$0.068} \\
 & Nov. & 0.354\,{\scriptsize$\pm$0.048} & 0.933\,{\scriptsize$\pm$0.317} & 0.250\,{\scriptsize$\pm$0.028} & 0.122\,{\scriptsize$\pm$0.019} & 0.217\,{\scriptsize$\pm$0.030} & $-$0.525\,{\scriptsize$\pm$0.054} & 0.226\,{\scriptsize$\pm$0.044} \\
\midrule
\multirow{2}{*}{w/o kernel} & May & 0.302\,{\scriptsize$\pm$0.054} & 1.017\,{\scriptsize$\pm$0.197} & 0.109\,{\scriptsize$\pm$0.052} & 0.911\,{\scriptsize$\pm$0.008} & 0.640\,{\scriptsize$\pm$0.072} & $-$0.466\,{\scriptsize$\pm$0.080} & 0.221\,{\scriptsize$\pm$0.055} \\
 & Nov. & 0.287\,{\scriptsize$\pm$0.060} & 0.890\,{\scriptsize$\pm$0.276} & \textbf{0.084}\,{\scriptsize$\pm$0.017} & 0.898\,{\scriptsize$\pm$0.021} & 0.574\,{\scriptsize$\pm$0.095} & $-$0.464\,{\scriptsize$\pm$0.043} & 0.211\,{\scriptsize$\pm$0.044} \\
\bottomrule
\end{tabular}
\caption{Ablation of Red-LGCP on May and Nov.\ 2025 (same windows and covariates as Table~\ref{tab:results-main}; mean $\pm$ population std over 10 seeds $\times$ 4 weekly windows per month). Each variant removes or replaces one component of the full model: ``w/o uncertainty'' uses the fixed threshold $\tau$ instead of $\tau+\sigma_{j,t}$; ``w/o redistribution'' zeroes rejected cells without recovering any of their expected counts; ``w/o gate'' reports the expected counts without zeroing; ``two-stage'' is not jointly optimized: the LGCP is fitted alone, with a Poisson likelihood, and the classifier is trained afterwards and used as a post-hoc gate with the fixed threshold $\tau$; ``w/o classifier'' is the LGCP alone, with a Poisson likelihood; ``w/o kernel'' removes the latent Gaussian process, keeping the covariate effects and the classifier. The first four rows share the same trained model, so their log-likelihood and CRPS, which score the hurdle predictive distribution, coincide; for ``two-stage'' and ``w/o classifier'' these scores refer to the predictive distribution of the Poisson LGCP, since a post-hoc gate only changes point predictions. Bold marks the best value per column within each month.}
\label{tab:results-ablation}
\end{table*}
